\documentclass[11pt]{article}
\usepackage{notes} % place notes.sty in the same folder
\definecolor{AccN}{HTML}{E05218}
\fancyhead[L]{\small\color{NavyN}\nouppercase{\leftmark}}\fancyhead[R]{\small\color{GreyN}Instructor Manual --- Lecture 1}
\usepackage{longtable,multicol,needspace}
\setlength{\emergencystretch}{2em}
\newtcolorbox{instr}[1]{nb,colback=AccN!6,colframe=AccN,title={#1}}
\newtcolorbox{prob}{nb,colback=BluN!5,colframe=BluN,title={Problem statement}}
\newtcolorbox{soln}{nb,colback=GrnN!5,colframe=GrnN,title={Model solution and marking points}}
\newtcolorbox{sayit}{nb,colback=white,colframe=NavyN,title={Read aloud to students}}
\newtcolorbox{live}[1]{nb,colback=AccN!4,colframe=AccN,title={#1}}
\begin{document}
\begin{titlepage}\centering\vspace*{2.2cm}
{\color{AccN}\large\bfseries LECTURE 1\par}\vspace{0.6cm}
{\color{NavyN}\Huge\bfseries Foundations, Paradigms,\\Defence Applications \& Data Pitfalls\par}\vspace{0.6cm}
{\color{AccN}\rule{0.7\textwidth}{1.5pt}}\par\vspace{0.6cm}
{\LARGE Instructor Manual\par}\vspace{0.4cm}
{\large\begin{minipage}{0.85\textwidth}\centering AI Foundations \textbullet\ ML Paradigms \textbullet\ Defence \& Wargaming Applications \textbullet\ Data Pitfalls\end{minipage}\par}\vspace{2cm}
{\large AK DHAMIJA, Researcher \& Instructor, AI \& Machine Learning\par}\vspace{0.2cm}
{\normalsize \href{mailto:dhamija.ak@gmail.com}{dhamija.ak@gmail.com} \quad\textbullet\quad 9811453163\par}\vspace{0.3cm}
{Unit 4: Simulation Strategies in Warfare --- AI \& Machine Learning in Warfare\par}\vfill
\begin{tcolorbox}[colback=NavyN!4,colframe=NavyN,boxrule=0.6pt,arc=2pt,width=0.9\textwidth]\small
\textbf{This manual contains:} the 45-minute session plan with timings; a slide-by-slide teaching guide for all 54 slides; key messages, facilitation tips, anticipated questions with model answers, misconceptions and in-class activities; live demonstrations with the interactive studio; consistency notes on the source material; the complete quiz answer key with rationale and scoring rules; and the assignment with marking scheme and full model solutions.
\end{tcolorbox}
\end{titlepage}
\tableofcontents\newpage


\needspace{12\baselineskip}\section{Lecture at a Glance}

\begin{center}\small\renewcommand{\arraystretch}{1.25}\begin{tabularx}{\textwidth}{@{}>{\bfseries}p{3.3cm}>{\raggedright\arraybackslash}X@{}}\toprule

\textcolor{NavyN}{Lecture}&1 --- Foundations, Paradigms, Defence Applications \& Data Pitfalls\\
\textcolor{NavyN}{Duration}&45 minutes: 5 min opening Q\&A, 30 min teaching, 10 min auto-quiz\\
\textcolor{NavyN}{Slides}&54 (Beamer PDF, 16:9), 17 sections\\
\textcolor{NavyN}{Prerequisite}&None (first lecture).\\
\textcolor{NavyN}{Student materials}&\texttt{Lecture1\_Slides\_Beamer.pdf}; \texttt{Lecture1\_Reference\_Text.pdf}; \texttt{Lecture1\_Quiz.html}; \texttt{Lecture1\_Answers\_Assignment.html} (release after the quiz)\\
\textcolor{NavyN}{Interactive studio}&\texttt{Lecture1\_Data\_Can\_Lie\_Studio.html} (live demonstrations in Section~\ref{sec:live})\\
\textcolor{NavyN}{Assessment}&Auto-quiz: 15 MCQs, 40~s each, $+1/-0.25$. Assignment: 5 problems, 100 marks.\\\bottomrule\end{tabularx}\end{center}

\subsection*{Overview}Lecture~1 sets the tone for the whole unit. Its first half (``Data Can Lie'') is deliberately placed \emph{before} any algorithm: students who later build classifiers, clusterers and RL agents must already distrust headline numbers, ask how data came to exist, and recompute claims from counts. The second half is a fast, wide survey of AI and machine learning --- classical AI, the paradigm landscape (supervised, unsupervised, semi- and self-supervised, reinforcement, deep and generative AI including LLMs and RAG) and the defence and wargaming application map --- that Lectures~2--8 then deepen one paradigm at a time. It closes by returning to data pitfalls (survivorship bias, overfitting, distribution shift, label noise) so that the lecture ends where it began: \emph{data and AI must be handled with care}.

\subsection*{Learning outcomes}By the end of the lecture students should be able to:\begin{enumerate}\item Explain, with counts, why Simpson's paradox arises and recompute an aggregate rate correctly.\item State what a $p$-value is and is not; distinguish statistical from practical significance; explain power and multiple testing.\item Explain why correlation does not imply causation (confounding, reverse causation, chance) and recognise regression to the mean.\item Apply the base-rate (Bayes) argument to compute the positive predictive value of an alarm.\item Describe the main AI/ML paradigms with one defence example each, including LLMs and RAG.\item Identify shortcut learning, survivorship bias, overfitting and distribution shift in a described system.\end{enumerate}

\subsection*{Syllabus mapping}\begin{itemize}\item AI \& ML paradigms --- supervised, unsupervised, reinforcement learning (overview; plus semi-/self-supervised, deep and generative AI)\item Foundations for all later syllabus items: data quality, statistical reasoning and interpretation of AI outputs\end{itemize}

\needspace{12\baselineskip}\section{Session Plan (45 Minutes)}

\begin{center}\begin{tikzpicture}[x=0.3444cm,y=1cm]
\fill[AmbN!40](0,0)rectangle(5,1);\node[font=\scriptsize\bfseries,align=center]at(2.5,0.5){Q\&A\\recap};
\fill[AccN!25](5.00,0)rectangle(6.00,1);\draw[white,thick](5.00,0)--(5.00,1);
\fill[NavyN!18](6.00,0)rectangle(14.50,1);\draw[white,thick](6.00,0)--(6.00,1);
\node[font=\tiny,rotate=90,anchor=center]at(10.25,0.5){\parbox{0.9cm}{\centering \S0}};
\fill[AccN!25](14.50,0)rectangle(15.50,1);\draw[white,thick](14.50,0)--(14.50,1);
\fill[NavyN!18](15.50,0)rectangle(16.00,1);\draw[white,thick](15.50,0)--(15.50,1);
\fill[AccN!25](16.00,0)rectangle(19.00,1);\draw[white,thick](16.00,0)--(16.00,1);
\node[font=\tiny,rotate=90,anchor=center]at(17.50,0.5){\parbox{0.9cm}{\centering \S3}};
\fill[NavyN!18](19.00,0)rectangle(19.50,1);\draw[white,thick](19.00,0)--(19.00,1);
\fill[AccN!25](19.50,0)rectangle(20.00,1);\draw[white,thick](19.50,0)--(19.50,1);
\fill[NavyN!18](20.00,0)rectangle(22.50,1);\draw[white,thick](20.00,0)--(20.00,1);
\node[font=\tiny,rotate=90,anchor=center]at(21.25,0.5){\parbox{0.9cm}{\centering \S6}};
\fill[AccN!25](22.50,0)rectangle(24.00,1);\draw[white,thick](22.50,0)--(22.50,1);
\node[font=\tiny,rotate=90,anchor=center]at(23.25,0.5){\parbox{0.9cm}{\centering \S7}};
\fill[NavyN!18](24.00,0)rectangle(25.00,1);\draw[white,thick](24.00,0)--(24.00,1);
\fill[AccN!25](25.00,0)rectangle(26.00,1);\draw[white,thick](25.00,0)--(25.00,1);
\fill[NavyN!18](26.00,0)rectangle(27.00,1);\draw[white,thick](26.00,0)--(26.00,1);
\fill[AccN!25](27.00,0)rectangle(28.50,1);\draw[white,thick](27.00,0)--(27.00,1);
\node[font=\tiny,rotate=90,anchor=center]at(27.75,0.5){\parbox{0.9cm}{\centering \S11}};
\fill[NavyN!18](28.50,0)rectangle(29.00,1);\draw[white,thick](28.50,0)--(28.50,1);
\fill[AccN!25](29.00,0)rectangle(30.00,1);\draw[white,thick](29.00,0)--(29.00,1);
\fill[NavyN!18](30.00,0)rectangle(34.00,1);\draw[white,thick](30.00,0)--(30.00,1);
\node[font=\tiny,rotate=90,anchor=center]at(32.00,0.5){\parbox{0.9cm}{\centering \S14}};
\fill[AccN!25](34.00,0)rectangle(35.00,1);\draw[white,thick](34.00,0)--(34.00,1);
\fill[RedN!30](35,0)rectangle(45,1);\node[font=\scriptsize\bfseries,align=center]at(40,0.5){Auto-quiz\\15 Q $\times$ 40 s};
\draw[NavyN](0,0)rectangle(45,1);
\foreach \tt in {0,5,10,15,20,25,30,35,40,45}{\draw[NavyN](\tt,0)--(\tt,-0.12);\node[font=\tiny,below]at(\tt,-0.12){\tt};}
\node[font=\scriptsize,below]at(22.5,-0.45){minutes};
\end{tikzpicture}\end{center}

\begin{center}\small\renewcommand{\arraystretch}{1.2}\begin{tabularx}{\textwidth}{@{}l l X r@{}}\toprule\textbf{Clock} & \textbf{Slides} & \textbf{Segment} & \textbf{Min}\\\midrule

0:00--5:00 & --- & \textbf{Opening Q\&A and recap} (see prompts below) & 5\\

5:00--6:00 & 2--3 & Course Overview (2 slides) & 1\\

6:00--14:30 & 4--16 & 0. Data Can Lie (13 slides) & 8.5\\

14:30--15:30 & 17--18 & 1. Intelligence \& AI (2 slides) & 1\\

15:30--16:00 & 19 & 2. History of AI (1 slide) & 0.5\\

16:00--19:00 & 20--24 & 3. Classical AI (5 slides) & 3\\

19:00--19:30 & 25 & 4. AI vs ML vs DL (1 slide) & 0.5\\

19:30--20:00 & 26 & 5. ML Paradigms (1 slide) & 0.5\\

20:00--22:30 & 27--30 & 6. Supervised Learning (4 slides) & 2.5\\

22:30--24:00 & 31--33 & 7. Unsupervised Learning (3 slides) & 1.5\\

24:00--25:00 & 34--35 & 8. Semi-sup \& SSL (2 slides) & 1\\

25:00--26:00 & 36--37 & 9. Reinforcement Learning (2 slides) & 1\\

26:00--27:00 & 38--39 & 10. Deep Learning (2 slides) & 1\\

27:00--28:30 & 40--42 & 11. Generative AI (3 slides) & 1.5\\

28:30--29:00 & 43 & 12. Defence Apps (1 slide) & 0.5\\

29:00--30:00 & 44--45 & 13. Wargaming Cases (2 slides) & 1\\

30:00--34:00 & 46--52 & 14. Data Pitfalls (7 slides) & 4\\

34:00--35:00 & 53--54 & 15. Summary (2 slides) & 1\\

35:00--45:00 & --- & \textbf{Auto-quiz} (instructions 1 min, 15 questions $\times$ 40 s = 10 min of answering; collect scores) & 10\\\bottomrule\end{tabularx}\end{center}

Slide 1 is the title slide and is shown while students settle. The teaching segment covers 53 content slides in 30 minutes, an average pace of 1.8 slides per minute. 

\begin{pitfall}{pace}This pace is only achievable if most survey slides get one sentence each. See the facilitation tips for a fast path through a single session and a two-session alternative.\end{pitfall}

\begin{instr}{Opening Q\&A prompts (first 5 minutes)}\begin{enumerate}\item (Icebreaker) Name one decision in your own work that was made on the basis of a number. How sure were you of that number?\item Where have you seen ``AI'' used in a defence or wargaming context? What did it actually do?\item What do you think is more dangerous: a bad algorithm or bad data? Hold that thought for 40 minutes.\end{enumerate}Take two or three answers per prompt; correct misconceptions briefly and note any topic that needs revisiting.\end{instr}

\needspace{12\baselineskip}\section{Slide-by-Slide Teaching Guide}

Each subsection gives the slide range, the suggested time, what to emphasise, and the slide titles in order. Numbers refer to pages of the Beamer PDF.

\needspace{8\baselineskip}\subsection{Course Overview \normalfont\small\color{GreyN}(slides 2--3, about 1 min)}

\textbf{Teaching note.} Show the 10-lecture roadmap and the lecture format (5-minute Q\&A, teaching, 10-minute quiz with $+1/-0.25$ marking). State plainly that Lecture~1 starts with statistics on purpose.

\begin{center}\small\renewcommand{\arraystretch}{1.1}\begin{tabularx}{\textwidth}{@{}r X@{}}2 & Course Overview --- 10-Lecture Roadmap (Unit 4: AI \& ML in Warfare)\\
3 & Lecture 1 --- Structure, Format \& Learning Objectives\\\end{tabularx}\end{center}

\needspace{8\baselineskip}\subsection{0. Data Can Lie \normalfont\small\color{GreyN}(slides 4--16, about 8.5 min)}

\textbf{Teaching note.} The heart of the lecture. Work Simpson's paradox \emph{on the board with counts}: A 81/87 and 192/263 versus B 234/270 and 55/80; totals 273/350 = 78\% versus 289/350 = 83\%. Then show the wargaming twin. Move briskly through the $p$-value, power (the 40-sortie coating trial, $p=0.072$, about 93 sorties per arm needed), multiple testing ($1-0.95^{20}=64\%$), Anscombe, spurious correlation (storks, $r\approx0.62$), regression to the mean (flight instructors), the base-rate fallacy (99\%/99\% detector at 1-in-1,000 gives PPV $\approx9\%$), Challenger, shortcut learning and automation bias. End on the Analyst's Oath slide and leave it up for a few seconds.

\begin{center}\small\renewcommand{\arraystretch}{1.1}\begin{tabularx}{\textwidth}{@{}r X@{}}4 & Before Any Algorithm: Data Can Lie --- Six Famous Failures, Six Disciplines\\
5 & Simpson's Paradox (I) --- Never Add or Average Percentages: Use the Counts\\
6 & Simpson's Paradox (II) --- The Wargaming Twin: ``AI-Assisted Units Perform Worse''?\\
7 & The Null Hypothesis \& Statistical Significance --- What a p-value Says\\
8 & Statistical vs Practical Significance --- and the p-hacking Trap\\
9 & Linear Correlation --- One Number Cannot Describe a Relationship\\
10 & Correlation $\neq$ Causation --- Spurious Correlation and How It Arises\\
11 & Regression to the Mean --- When Chance Masquerades as Cause\\
12 & The Base-Rate Fallacy --- Why a ``99\% Accurate'' Detector Mostly Cries Wolf\\
13 & Engineering Lesson: Challenger (1986) --- The Data You Leave Out\\
14 & AI Interpretation Errors (I) --- Shortcut Learning: Right Answer, Wrong Reason\\
15 & AI Interpretation Errors (II) --- When Humans Misread the Machine\\
16 & The Analyst's Oath --- Ten Rules Before Trusting Any Data or AI Output\\\end{tabularx}\end{center}

\needspace{8\baselineskip}\subsection{1. Intelligence \& AI \normalfont\small\color{GreyN}(slides 17--18, about 1 min)}

\textbf{Teaching note.} Define intelligence through its four cognitive pillars and AI through the classic definitions and the Turing test. Keep it short; the point is vocabulary.

\begin{center}\small\renewcommand{\arraystretch}{1.1}\begin{tabularx}{\textwidth}{@{}r X@{}}17 & What is Intelligence? The Four Cognitive Pillars\\
18 & What is Artificial Intelligence? Definitions and the Turing Test\\\end{tabularx}\end{center}

\needspace{8\baselineskip}\subsection{2. History of AI \normalfont\small\color{GreyN}(slide 19, about 0.5 min)}

\textbf{Teaching note.} One slide of history: symbolic AI, the AI winters, the statistical turn, deep learning (2012) and foundation models. Emphasise that each boom was followed by over-promising --- a recurring defence-procurement risk.

\begin{center}\small\renewcommand{\arraystretch}{1.1}\begin{tabularx}{\textwidth}{@{}r X@{}}19 & Brief History of Artificial Intelligence\\\end{tabularx}\end{center}

\needspace{8\baselineskip}\subsection{3. Classical AI \normalfont\small\color{GreyN}(slides 20--24, about 3 min)}

\textbf{Teaching note.} Expert systems and search (game trees, A*): they are transparent and still used, but brittle and they suffer from state-space explosion. Contrast the numerical results with the brittleness slide; ask why a rule base fails when the enemy changes tactics.

\begin{center}\small\renewcommand{\arraystretch}{1.1}\begin{tabularx}{\textwidth}{@{}r X@{}}20 & Classical AI --- Expert Systems in Defence\\
21 & Classical AI --- Expert Systems: Numerical Performance \& Brittleness\\
22 & Classical AI --- Search Algorithms: Game Trees \& A*\\
23 & Classical AI --- A* Planning and the State-Space Explosion\\
24 & Why Classical AI Fails in Real Warfare\\\end{tabularx}\end{center}

\needspace{8\baselineskip}\subsection{4. AI vs ML vs DL \normalfont\small\color{GreyN}(slide 25, about 0.5 min)}

\textbf{Teaching note.} The nested picture AI $\supset$ ML $\supset$ DL. Clarify that not all AI learns and not all ML is deep.

\begin{center}\small\renewcommand{\arraystretch}{1.1}\begin{tabularx}{\textwidth}{@{}r X@{}}25 & AI vs Machine Learning vs Deep Learning --- The Nested Hierarchy\\\end{tabularx}\end{center}

\needspace{8\baselineskip}\subsection{5. ML Paradigms \normalfont\small\color{GreyN}(slide 26, about 0.5 min)}

\textbf{Teaching note.} The ``Big Five'' map. This slide is the index for the whole unit: point to where each paradigm is taught (L2--3, L4--5, L6--8).

\begin{center}\small\renewcommand{\arraystretch}{1.1}\begin{tabularx}{\textwidth}{@{}r X@{}}26 & Machine Learning Paradigms --- The Big Five\\\end{tabularx}\end{center}

\needspace{8\baselineskip}\subsection{6. Supervised Learning \normalfont\small\color{GreyN}(slides 27--30, about 2.5 min)}

\textbf{Teaching note.} Supervised learning concept, evaluation protocol and the radar case study. Recompute the confusion-matrix diagonal with the class: $493+488+487+489+496=2{,}453$ of 2,500, i.e.\ 98.1\% (see the consistency notes). Stress the asymmetric cost of a false negative.

\begin{center}\small\renewcommand{\arraystretch}{1.1}\begin{tabularx}{\textwidth}{@{}r X@{}}27 & Supervised Learning --- Concept \& Mathematics\\
28 & Supervised Learning --- Evaluation Protocol \& Key Metrics\\
29 & Supervised Learning --- Radar Signal Classification (Case Study)\\
30 & Supervised Learning --- The CNN Confusion Matrix and the Cost of a False Negative\\\end{tabularx}\end{center}

\needspace{8\baselineskip}\subsection{7. Unsupervised Learning \normalfont\small\color{GreyN}(slides 31--33, about 1.5 min)}

\textbf{Teaching note.} Unsupervised learning, the elbow and silhouette, and OPFOR tactic clustering. Preview that Lecture~4 discovers an archetype nobody had labelled.

\begin{center}\small\renewcommand{\arraystretch}{1.1}\begin{tabularx}{\textwidth}{@{}r X@{}}31 & Unsupervised Learning --- Concept \& Core Algorithms\\
32 & Unsupervised Learning --- Choosing $k$: the Elbow Method \& Silhouette\\
33 & Unsupervised Learning --- OPFOR Tactic Clustering (Case Study)\\\end{tabularx}\end{center}

\needspace{8\baselineskip}\subsection{8. Semi-sup \& SSL \normalfont\small\color{GreyN}(slides 34--35, about 1 min)}

\textbf{Teaching note.} Semi-supervised and self-supervised learning: why defence has abundant unlabelled data and scarce labels, and how pretext tasks create labels from the data itself.

\begin{center}\small\renewcommand{\arraystretch}{1.1}\begin{tabularx}{\textwidth}{@{}r X@{}}34 & Semi-supervised Learning --- Learning with Limited Labels\\
35 & Self-supervised Learning --- Labels from the Data Itself\\\end{tabularx}\end{center}

\needspace{8\baselineskip}\subsection{9. Reinforcement Learning \normalfont\small\color{GreyN}(slides 36--37, about 1 min)}

\textbf{Teaching note.} Reinforcement learning as an MDP; the Bellman equation, Q-learning and DQN at headline level only. Lectures~6--8 do the mathematics.

\begin{center}\small\renewcommand{\arraystretch}{1.1}\begin{tabularx}{\textwidth}{@{}r X@{}}36 & Reinforcement Learning --- Concept \& Markov Decision Process\\
37 & Reinforcement Learning --- Bellman Equation, Q-Learning and DQN\\\end{tabularx}\end{center}

\needspace{8\baselineskip}\subsection{10. Deep Learning \normalfont\small\color{GreyN}(slides 38--39, about 1 min)}

\textbf{Teaching note.} The artificial neuron, depth as hierarchical features, and the main architectures (CNN, RNN/LSTM, Transformer, GAN) with defence data types.

\begin{center}\small\renewcommand{\arraystretch}{1.1}\begin{tabularx}{\textwidth}{@{}r X@{}}38 & Deep Learning --- Neural Networks \& the Neuron\\
39 & Deep Learning --- Why ``Deep''? Hierarchical Features \& Architectures\\\end{tabularx}\end{center}

\needspace{8\baselineskip}\subsection{11. Generative AI \normalfont\small\color{GreyN}(slides 40--42, about 1.5 min)}

\textbf{Teaching note.} How LLMs work (next-token prediction), their defence uses and risks (hallucination, leakage, prompt injection), and RAG as a way to ground answers in an approved document store. Demonstrate a hallucinated citation if you can.

\begin{center}\small\renewcommand{\arraystretch}{1.1}\begin{tabularx}{\textwidth}{@{}r X@{}}40 & Generative AI --- How Large Language Models Work\\
41 & Generative AI --- LLMs in Defence: Applications and Risks\\
42 & Retrieval Augmented Generation (RAG) for Defence Intelligence\\\end{tabularx}\end{center}

\needspace{8\baselineskip}\subsection{12. Defence Apps \normalfont\small\color{GreyN}(slide 43, about 0.5 min)}

\textbf{Teaching note.} The application landscape: ISR, EW, cyber, logistics, C2 decision support, training and simulation.

\begin{center}\small\renewcommand{\arraystretch}{1.1}\begin{tabularx}{\textwidth}{@{}r X@{}}43 & AI/ML in Defence --- Application Landscape\\\end{tabularx}\end{center}

\needspace{8\baselineskip}\subsection{13. Wargaming Cases \normalfont\small\color{GreyN}(slides 44--45, about 1 min)}

\textbf{Teaching note.} Computer-generated forces, attrition modelling, terrain analysis, COA generation and maturity levels. Ask the class where in a wargame AI would help \emph{and} where it would mislead.

\begin{center}\small\renewcommand{\arraystretch}{1.1}\begin{tabularx}{\textwidth}{@{}r X@{}}44 & AI in Wargaming --- Computer Generated Forces and Attrition Modelling\\
45 & AI in Wargaming --- Terrain, COA Generation and Maturity Levels\\\end{tabularx}\end{center}

\needspace{8\baselineskip}\subsection{14. Data Pitfalls \normalfont\small\color{GreyN}(slides 46--52, about 4 min)}

\textbf{Teaching note.} Return to data pitfalls in the ML setting: survivorship bias, spurious correlation, overfitting (mine detectors), distribution shift (KS statistic; 94.1\% Afghanistan to 58.7\% Iraq), training-data bias and label noise (97.2\% falls to 61.7\% with 30\% label noise).

\begin{center}\small\renewcommand{\arraystretch}{1.1}\begin{tabularx}{\textwidth}{@{}r X@{}}46 & Data Pitfalls --- Why Bad Data Produces Dangerous AI\\
47 & Data Pitfall \#1 --- Survivorship Bias: Studying Only the Winners\\
48 & Data Pitfall \#2 --- Spurious Correlation \& Simpson's Paradox\\
49 & Data Pitfall \#2 (cont.) --- Misuse of Statistics \& Shortcut Learning\\
50 & Data Pitfall \#3 --- Overfitting: the Mine Detection Systems\\
51 & Data Pitfall \#3 --- Distribution Shift: Detect It Before You Deploy\\
52 & Data Pitfalls \#4 \& \#5 --- Training Data Bias and Label Noise\\\end{tabularx}\end{center}

\needspace{8\baselineskip}\subsection{15. Summary \normalfont\small\color{GreyN}(slides 53--54, about 1 min)}

\textbf{Teaching note.} Summary and look-ahead. Re-read the three to five rules from the Analyst's Oath that you want students to remember.

\begin{center}\small\renewcommand{\arraystretch}{1.1}\begin{tabularx}{\textwidth}{@{}r X@{}}53 & Lecture 1 Summary --- Key Takeaways\\
54 & Looking Ahead --- Lecture 2: Supervised Learning (Classification)\\\end{tabularx}\end{center}

\needspace{12\baselineskip}\section{Key Messages}Students should leave the room able to state these in their own words:\begin{keypoint}\begin{enumerate}\item Never add or average percentages: totals come from counts, and a confounder can reverse a comparison.\item A $p$-value is a statement about the data assuming $H_0$, not the probability that $H_0$ is true; absence of evidence is not evidence of absence.\item Correlation is not causation; ask about confounders, reverse causation, selection and chance.\item Rare events make even accurate detectors produce mostly false alarms; always compute the positive predictive value.\item AI automates and scales whatever is in the data, including bias; models can be right for the wrong reason.\item Each ML paradigm answers a different question: labels (supervised), structure (unsupervised), decisions (RL), generation (GenAI).\end{enumerate}\end{keypoint}

\needspace{12\baselineskip}\section{Facilitation Tips}\begin{itemize}\item With 54 slides this is the densest lecture. In a single 45-minute session, teach the 13 ``Data Can Lie'' slides fully (about 14 minutes), and give Sections 1--13 one sentence per slide, pointing students to the reference text for depth. If the timetable allows, run it as two sessions: \textbf{1a} Course overview + Data Can Lie + Analyst's Oath; \textbf{1b} AI paradigms to the end.\item Do the Simpson arithmetic live. Ask the class to predict the overall winner \emph{before} you reveal the totals; most will say A. That surprise is the lesson.\item For the base-rate slide, use natural frequencies: ``out of 100,000 tracks, 100 are threats; 99 are caught; 999 of the 99,900 non-threats also alarm''.\item Keep the AI-survey slides moving. Students will meet each paradigm again in depth; today they only need the map.\item Leave the Analyst's Oath visible while you take the last questions before the quiz.\end{itemize}

\needspace{12\baselineskip}\section{Anticipated Questions and Model Answers}

\begin{question}{anticipated question}\textbf{Q:} Isn't the kidney-stone example just medicine? Why is it in a defence course?\par\smallskip\textbf{A:} Because the structure is identical whenever assignment to a ``treatment'' depends on difficulty: AI tools sent to the hardest sectors, new tactics tried in the worst terrain, experienced crews given the toughest missions. The wargaming-twin slide makes the defence version explicit.\end{question}

\begin{question}{anticipated question}\textbf{Q:} If $p=0.072$, isn't the coating ``almost significant''?\par\smallskip\textbf{A:} The trial was under-powered: with 40 sorties per arm the 95\% CI for the difference runs from about $-1$ to $+41$ points. The right conclusion is ``inconclusive --- run about 93 sorties per arm'', not ``no effect'' and not ``almost works''.\end{question}

\begin{question}{anticipated question}\textbf{Q:} Can't a big enough deep network learn its way around biased data?\par\smallskip\textbf{A:} No. A network fits the data distribution it is given. If the data are biased or the deployment distribution differs, more capacity can make the model fit the bias more precisely. The Afghanistan-to-Iraq mine-detector drop is the example.\end{question}

\begin{question}{anticipated question}\textbf{Q:} Is an LLM with RAG safe to use for intelligence summaries?\par\smallskip\textbf{A:} Safer, not safe. RAG grounds answers in retrieved documents and allows citations, but the model can still misread or over-generalise, and the document store itself must be curated and access-controlled. A human must check claims against the cited sources.\end{question}

\begin{question}{anticipated question}\textbf{Q:} Why does regression to the mean not appear in the textbook formulas?\par\smallskip\textbf{A:} It does: the regression slope of a re-test on a test is $r\,\sigma_y/\sigma_x$, which is below 1 when $r<1$. Extreme scores are partly luck, and luck does not repeat.\end{question}

\needspace{12\baselineskip}\section{Common Misconceptions}\begin{center}\small\renewcommand{\arraystretch}{1.35}\begin{tabularx}{\textwidth}{@{}>{\raggedright\arraybackslash}X >{\raggedright\arraybackslash}X@{}}\toprule\textbf{Misconception} & \textbf{Correction}\\\midrule``A $p$-value of 0.03 means a 3\% chance the null hypothesis is true.'' & It is $P(\text{data this extreme}\mid H_0)$, not $P(H_0\mid\text{data})$.\\
``Overall rates can be found by averaging subgroup rates.'' & Only if the subgroups are equal in size; otherwise weight by counts.\\
``A 99\%-accurate detector means 99\% of its alarms are real.'' & PPV depends on the base rate; at 1 in 1,000 it is about 9\%.\\
``Strong correlation proves a relationship worth acting on.'' & It may be confounded, reversed or spurious; it may also be non-linear (Anscombe).\\
``Deep learning has replaced classical AI.'' & Search, planning and rules remain essential where transparency and guarantees matter; hybrids are common.\\\bottomrule\end{tabularx}\end{center}

\needspace{12\baselineskip}\section{In-Class Activities and Discussion Prompts}\label{sec:act}\begin{itemize}\item \textbf{Predict-then-reveal (3 min):} show only the subgroup rates of the wargaming Simpson table and ask for a show of hands on which unit type is better overall.\item \textbf{Spot the flaw (5 min):} read out a one-paragraph ``after-action report'' claiming that units with more maintenance had more breakdowns, so maintenance should be cut. Pairs name the error (reverse causation) and propose the data needed.\item \textbf{Base-rate ladder:} ask what PPV becomes if prevalence rises from 1 in 1,000 to 1 in 100 (about 50\%). Discuss what intelligence cueing does to the base rate.\end{itemize}

\needspace{12\baselineskip}\section{Run It Live: Studio Demonstrations}\label{sec:live}

Project \texttt{Lecture1\_Data\_Can\_Lie\_Studio.html} (it is a single HTML file and runs offline in any browser; every result updates as an input changes). Each demonstration below follows the same rhythm: \textbf{set} the inputs, \textbf{ask} the class to commit to a prediction (a show of hands or a one-word answer), then \textbf{reveal} the result and \textbf{land} the point. The numbers are exactly what the page shows for the settings given; simulations use fixed random seeds, so they repeat every time. The 45-minute plan has little slack: choose the one demonstration marked \emph{first choice} for class, and use the others in a tutorial, a second session, or as homework with the drill cards at the bottom of the page.

\begin{live}{Demo 1: Why a 99\% detector mostly cries wolf (about 3 min, with \S0 Data Can Lie or \S14 Data Pitfalls) --- \emph{first choice}}\textbf{Set.} Trap~5 at its defaults: sensitivity 0.99, specificity 0.99, 1 threat per 1,000 tracks.\par\smallskip\textbf{Ask.} ``Of the alarms this detector raises, what share are real threats: over 90\%, about half, or under 20\%?''\par\smallskip\textbf{Reveal.} Only \textbf{9.0\%}: per 100,000 tracks there are 99 true alarms and 999 false ones. Now type 50 threats per 1,000: precision jumps to \textbf{83.9\%} (4,950 true, 950 false).\par\smallskip\textbf{Land the point.} The detector did not change; the population did. Precision belongs to a detector \emph{and} the base rate it meets, which is why a second, independent sensor is used to confirm.\end{live}

\begin{live}{Demo 2: Test enough tactics and one will ``work'' (about 2 min, with \S14 Data Pitfalls)}\textbf{Set.} Trap~3: 20 tactics tested, $\alpha=0.05$, and none of them truly works.\par\smallskip\textbf{Ask.} ``How likely is it that at least one tactic comes out `significant'?''\par\smallskip\textbf{Reveal.} \textbf{64\%}, with 1.00 false discovery expected; the Bonferroni threshold per test is 0.0025. Drag to 100 tactics: \textbf{99\%}, 5.00 false discoveries, per-test threshold 0.0005.\par\smallskip\textbf{Land the point.} A single significant result from a large search is weak evidence. Fix the hypothesis in advance, correct for multiple tests, and replicate before briefing.\end{live}

\begin{live}{Demo 3: The flight-instructor illusion (about 2 min, with \S0 Data Can Lie)}\textbf{Set.} Trap~4: test--retest correlation $\rho=0.5$, top and bottom 10\% selected. The model contains \emph{no} effect of praise or criticism.\par\smallskip\textbf{Ask.} ``Instructors praise the best landings and criticise the worst. Next time the praised group does worse and the criticised group better. Does criticism work?''\par\smallskip\textbf{Reveal.} Praised group 87.5 $\rightarrow$ \textbf{78.8}; criticised group 52.5 $\rightarrow$ \textbf{61.2}, from chance alone. Raise $\rho$ to 0.9: 87.5 $\rightarrow$ 85.8 and 52.5 $\rightarrow$ 54.2.\par\smallskip\textbf{Land the point.} Extreme results regress toward the mean, more so when the measure is unreliable. Only a comparison group can show that an intervention worked.\end{live}

\needspace{12\baselineskip}\section{Consistency Notes on the Source Material}These notes record where the original PowerPoint deck contained an inconsistency and how the current slides, reference text, quiz and assignment handle it. Students who have seen the old deck may raise them.\begin{itemize}\item The original PPTX Simpson slide claimed an overall tie (83\% vs 83\%) while its own counts gave 166 vs 156. All new materials use the published kidney-stone counts: A 273/350 = 78\%, B 289/350 = 83\%; the corrected PPTX shows counts alongside percentages.\item The radar confusion matrix has a diagonal of 2,453/2,500 = 98.1\%, while the headline accuracy is 97.2\%; the text mentions three bombers labelled friendly where the matrix shows two. The slides keep the source figures; use the discrepancy as a live ``recompute the headline'' exercise.\end{itemize}

\needspace{12\baselineskip}\section{Quiz Administration and Answer Key}

\begin{sayit}``You will now take a 15-question quiz. Each question has four options and exactly one correct answer. You have 40 seconds per question; the quiz advances automatically when time runs out. A correct answer scores $+1$, a wrong answer $-0.25$, and an unanswered question 0. The maximum is 15 and the minimum is $-3.75$. Because of the negative marking, guessing blindly among four options gains nothing on average --- answer when you can eliminate options. Your score and a full review appear at the end.''\end{sayit}

\textbf{Delivery.} Open \texttt{Lecture1\_Quiz.html} on each student device (or project it and collect answers on paper). The page enforces the timer and computes the score with negative marking. Release \texttt{Lecture1\_Answers\_Assignment.html} only after all students have finished.

\subsection*{Answer key}\begin{center}\small\renewcommand{\arraystretch}{1.3}\begin{tabular}{|c|c|c|c|c|c|c|c|c|c|c|c|c|c|c|}\hline \textbf{1} & \textbf{2} & \textbf{3} & \textbf{4} & \textbf{5} & \textbf{6} & \textbf{7} & \textbf{8} & \textbf{9} & \textbf{10} & \textbf{11} & \textbf{12} & \textbf{13} & \textbf{14} & \textbf{15}\\\hline D & A & D & B & C & A & C & B & B & C & C & C & A & D & D\\\hline\end{tabular}\end{center}

\textbf{Interpreting scores.} 12 or more: secure command of the lecture. 9--11.75: good, review missed items. Below 9: revisit the reference text before the assignment. If more than half the class misses a question, revisit that point in the next lecture's opening Q\&A.

\subsection*{Questions, options and rationale}

{\small

\noindent\textbf{Q1.} In the kidney-stone study, treatment A succeeded in 81/87 small-stone and 192/263 large-stone cases. What is A's overall success rate?\begin{itemize}[leftmargin=2.2em,itemsep=0pt,topsep=1pt]\item[A.] 83\%, the average of 93\% and 73\%\item[B.] 93\%, because A wins in every subgroup\item[C.] 166\%, the sum of the subgroup rates\item[\textbf{D.}] \textbf{78\%}~{\color{GrnN}$\checkmark$}\end{itemize}\noindent{\color{GreyN}\textit{Rationale:}} Totals come from counts: (81+192)/(87+263) = 273/350 = 78\%. Averaging or adding subgroup percentages ignores group sizes and is meaningless.\par\medskip

\noindent\textbf{Q2.} Treatment A beats B for both small and large stones, yet loses overall (78\% vs 83\%). The best explanation is:\begin{itemize}[leftmargin=2.2em,itemsep=0pt,topsep=1pt]\item[\textbf{A.}] \textbf{A was given mostly the harder large-stone cases, so its total is dragged down by case mix}~{\color{GrnN}$\checkmark$}\item[B.] Random chance\item[C.] B is genuinely the better treatment\item[D.] The data contain an arithmetic error\end{itemize}\noindent{\color{GreyN}\textit{Rationale:}} Simpson's paradox: stone size is a confounder affecting both treatment choice and outcome. 75\% of A's patients had large stones vs 23\% of B's.\par\medskip

\noindent\textbf{Q3.} A p-value of 0.03 means:\begin{itemize}[leftmargin=2.2em,itemsep=0pt,topsep=1pt]\item[A.] There is a 3\% chance that the null hypothesis is true\item[B.] There is a 97\% chance the effect is real\item[C.] The effect is small\item[\textbf{D.}] \textbf{If the null hypothesis were true, data at least this extreme would occur about 3\% of the time}~{\color{GrnN}$\checkmark$}\end{itemize}\noindent{\color{GreyN}\textit{Rationale:}} A p-value is computed assuming H0 is true. It is not the probability that H0 is true, and says nothing about effect size.\par\medskip

\noindent\textbf{Q4.} A coating trial with 40 sorties per arm finds detection 65\% vs 45\% (p = 0.072). The correct conclusion is:\begin{itemize}[leftmargin=2.2em,itemsep=0pt,topsep=1pt]\item[A.] The coating has no effect\item[\textbf{B.}] \textbf{The result is inconclusive: the trial was under-powered for an effect of this size}~{\color{GrnN}$\checkmark$}\item[C.] The coating reduces detection by exactly 20 points\item[D.] The coating increases detection\end{itemize}\noindent{\color{GreyN}\textit{Rationale:}} Absence of evidence is not evidence of absence. The 95\% CI for the difference is about -1 to +41 points; about 93 sorties per arm were needed for 80\% power.\par\medskip

\noindent\textbf{Q5.} Twenty independent tactics are each tested at alpha = 0.05 and none truly has an effect. The probability that at least one appears significant is about:\begin{itemize}[leftmargin=2.2em,itemsep=0pt,topsep=1pt]\item[A.] 5\%\item[B.] 20\%\item[\textbf{C.}] \textbf{64\%}~{\color{GrnN}$\checkmark$}\item[D.] 100\%\end{itemize}\noindent{\color{GreyN}\textit{Rationale:}} $1-(1-0.05)^{20}=1-0.95^{20}\approx0.64$. Correct with Bonferroni (alpha/m) and confirm on fresh data.\par\medskip

\noindent\textbf{Q6.} One million simulated engagements per tactic give hit rates 50.0\% vs 50.2\% with p = 0.005. This shows:\begin{itemize}[leftmargin=2.2em,itemsep=0pt,topsep=1pt]\item[\textbf{A.}] \textbf{A statistically significant but operationally trivial difference}~{\color{GrnN}$\checkmark$}\item[B.] A large, important improvement\item[C.] An invalid test\item[D.] That more data always reduces bias\end{itemize}\noindent{\color{GreyN}\textit{Rationale:}} With huge samples tiny effects become significant. Report effect sizes with confidence intervals and judge operational relevance.\par\medskip

\noindent\textbf{Q7.} Anscombe's quartet teaches that:\begin{itemize}[leftmargin=2.2em,itemsep=0pt,topsep=1pt]\item[A.] Correlation always implies causation\item[B.] Outliers never affect correlation\item[\textbf{C.}] \textbf{Datasets with identical means, variances, correlation (0.816) and regression line can look completely different, so always plot the data}~{\color{GrnN}$\checkmark$}\item[D.] r = 0 means no relationship\end{itemize}\noindent{\color{GreyN}\textit{Rationale:}} Summary statistics hide shape: a curve, an outlier or a single leverage point can produce the same r.\par\medskip

\noindent\textbf{Q8.} Units with the most maintenance hours have the most breakdowns. The most likely explanation is:\begin{itemize}[leftmargin=2.2em,itemsep=0pt,topsep=1pt]\item[A.] Maintenance causes breakdowns\item[\textbf{B.}] \textbf{Reverse causation: failing vehicles attract maintenance}~{\color{GrnN}$\checkmark$}\item[C.] A confounder: unit budget\item[D.] Simpson's paradox\end{itemize}\noindent{\color{GreyN}\textit{Rationale:}} The arrow runs from breakdowns to maintenance, not the other way.\par\medskip

\noindent\textbf{Q9.} Across 17 European countries stork pairs correlate with birth rates (r about 0.62). This is best explained by:\begin{itemize}[leftmargin=2.2em,itemsep=0pt,topsep=1pt]\item[A.] Storks deliver babies\item[\textbf{B.}] \textbf{A confounder such as land area or rural population}~{\color{GrnN}$\checkmark$}\item[C.] Measurement error\item[D.] Regression to the mean\end{itemize}\noindent{\color{GreyN}\textit{Rationale:}} A common cause drives both variables; the correlation is spurious.\par\medskip

\noindent\textbf{Q10.} Flight instructors notice that cadets praised after excellent landings usually do worse next time. This is mainly:\begin{itemize}[leftmargin=2.2em,itemsep=0pt,topsep=1pt]\item[A.] Evidence that praise harms performance\item[B.] Survivorship bias\item[\textbf{C.}] \textbf{Regression to the mean}~{\color{GrnN}$\checkmark$}\item[D.] The base-rate fallacy\end{itemize}\noindent{\color{GreyN}\textit{Rationale:}} An exceptional landing contains luck; the next one tends to be closer to the cadet's average regardless of feedback. Use a control group.\par\medskip

\noindent\textbf{Q11.} A detector has 99\% sensitivity and 99\% specificity; 1 in 1,000 tracks is a real threat. Roughly what fraction of alarms are real threats?\begin{itemize}[leftmargin=2.2em,itemsep=0pt,topsep=1pt]\item[A.] 99\%\item[B.] 50\%\item[\textbf{C.}] \textbf{9\%}~{\color{GrnN}$\checkmark$}\item[D.] 1\%\end{itemize}\noindent{\color{GreyN}\textit{Rationale:}} PPV = (0.99 x 0.001) / (0.99 x 0.001 + 0.01 x 0.999) is about 0.09. About 10 true alarms and 100 false ones per 10,000 tracks.\par\medskip

\noindent\textbf{Q12.} Wald recommended armouring returning bombers where:\begin{itemize}[leftmargin=2.2em,itemsep=0pt,topsep=1pt]\item[A.] Bullet holes were densest\item[B.] Uniformly over the whole aircraft\item[\textbf{C.}] \textbf{No bullet holes were seen, because aircraft hit there did not return}~{\color{GrnN}$\checkmark$}\item[D.] Only on the wings\end{itemize}\noindent{\color{GreyN}\textit{Rationale:}} Survivorship bias: the missing aircraft were hit in the engines and cockpit. The absence of data is the signal.\par\medskip

\noindent\textbf{Q13.} What was the key analytical error in the Challenger O-ring discussion?\begin{itemize}[leftmargin=2.2em,itemsep=0pt,topsep=1pt]\item[\textbf{A.}] \textbf{Examining mainly flights with damage, which hid the temperature trend shown by damage-free flights}~{\color{GrnN}$\checkmark$}\item[B.] Using too many data points\item[C.] Using a logistic model\item[D.] Ignoring the shuttle's weight\end{itemize}\noindent{\color{GreyN}\textit{Rationale:}} Excluding no-event cases removes the contrast that reveals the trend; the launch also extrapolated far below any tested temperature.\par\medskip

\noindent\textbf{Q14.} A classifier separates wolves from huskies by detecting snow in the background. This is an example of:\begin{itemize}[leftmargin=2.2em,itemsep=0pt,topsep=1pt]\item[A.] Overfitting to label noise\item[B.] Underfitting\item[C.] Class imbalance\item[\textbf{D.}] \textbf{Shortcut learning: right answer for the wrong reason}~{\color{GrnN}$\checkmark$}\end{itemize}\noindent{\color{GreyN}\textit{Rationale:}} The model exploited a spurious cue correlated with the label (as Clever Hans did); explanation tools such as LIME exposed it.\par\medskip

\noindent\textbf{Q15.} A mine detector scores 94.1\% in Afghanistan but 58.7\% in Iraq. The best pre-deployment safeguard is:\begin{itemize}[leftmargin=2.2em,itemsep=0pt,topsep=1pt]\item[A.] Train a bigger network on more Afghan data\item[B.] Increase the decision threshold\item[C.] Report training accuracy only\item[\textbf{D.}] \textbf{Test on data from the deployment domain and monitor feature distributions (e.g. the KS statistic)}~{\color{GrnN}$\checkmark$}\end{itemize}\noindent{\color{GreyN}\textit{Rationale:}} Overfitting to the training domain plus distribution shift. More biased data amplifies bias; only deployment-domain testing reveals the gap.\par\medskip

}

\needspace{12\baselineskip}\section{Assignment: Marking Scheme and Model Solutions}

\textbf{Instructions.} Attempt all five problems (100 marks, about 3 hours). Show every step of your working: marks are awarded for method as well as answers. Use the reference text for Lecture 1, Sections 3--11. Solutions are hidden --- attempt each problem before expanding its solution.

\begin{center}\small\renewcommand{\arraystretch}{1.25}\begin{tabularx}{0.85\textwidth}{@{}l X r@{}}\toprule\textbf{\#} & \textbf{Problem} & \textbf{Marks}\\\midrule1 & Simpson's paradox in COA-tool evaluation & 20\\
2 & Significance, power and multiple testing & 20\\
3 & Base rates and two-stage screening for insider threats & 20\\
4 & Regression to the mean and confounding & 20\\
5 & Distribution shift and evaluation design & 20\\\midrule & \textbf{Total} & \textbf{100}\\\bottomrule\end{tabularx}\end{center}

\begin{instr}{General marking guidance}\begin{itemize}
\item Marks for each part are shown in square brackets in the problem statement.
\item Award method marks for correct working even when an arithmetic slip carries through; deduct once for the slip, not again for its consequences.
\item For interpretive parts, full marks need (i) the correct conclusion, (ii) the reason, and (iii) an operational implication or caveat.
\item Answers that state a number without uncertainty, baseline or caveat where one is required lose a mark per part.
\item Suggested deadline: one week after the lecture; return marked work before the next quiz on the same topic.
\end{itemize}\end{instr}

\needspace{8\baselineskip}\subsection{Problem 1: Simpson's paradox in COA-tool evaluation \normalfont\small\color{GreyN}(20 marks)}

\begin{prob}Two course-of-action planning tools, X and Y, were used in 600 exercise engagements. Mission successes by threat level:

\begin{center}\small\begin{tabular}{@{}cccccccccccc@{}}\toprule Threat level&Tool X&Tool Y\\
Low&45 / 50&170 / 200\\
Medium&60 / 80&55 / 80\\
High&70 / 170&6 / 20\\\bottomrule
\end{tabular}\end{center}

(a) Compute each tool's success rate in each threat level and overall. [6]\newline (b) Explain precisely why the overall comparison misleads. [4]\newline (c) Compute \textit{directly standardised} success rates using the combined threat mix of all 600 engagements as the standard population. [6]\newline (d) Write a two-sentence recommendation to the programme manager. [4]\end{prob}

\begin{soln}\textbf{(a)} X: low 90.0\%, medium 75.0\%, high 41.2\%; overall 175/300 = \textbf{58.3\%}. Y: low 85.0\%, medium 68.75\%, high 30.0\%; overall 231/300 = \textbf{77.0\%}. X is better in every stratum (+5.0, +6.25, +11.2 points) yet 18.7 points worse overall.

\textbf{(b)} Threat level is a confounder: it affects success and was distributed very differently. 170/300 = 57\% of X's engagements were high-threat, against 20/300 = 7\% of Y's. The overall rate is a weighted average with weights equal to each tool's own threat mix, so X is penalised for its harder assignments (Simpson's paradox).

\textbf{(c)} Combined mix: low 250, medium 160, high 190 (total 600), so weights 0.4167, 0.2667, 0.3167. X: 0.9(0.4167) + 0.75(0.2667) + 0.4118(0.3167) = 0.375 + 0.200 + 0.130 = \textbf{70.5\%}. Y: 0.85(0.4167) + 0.6875(0.2667) + 0.30(0.3167) = 0.354 + 0.183 + 0.095 = \textbf{63.3\%}. On a common case mix X outperforms Y by about 7 points.

\textbf{(d)} Tool X performs better than Y at every threat level and by about 7 points on a standardised mix; its lower raw rate reflects its assignment to high-threat sectors. Retain X, and in future trials randomise tool assignment across threat levels so that raw comparisons are fair.\end{soln}

\needspace{8\baselineskip}\subsection{Problem 2: Significance, power and multiple testing \normalfont\small\color{GreyN}(20 marks)}

\begin{prob}A new jamming pod is trialled. Hostile radar achieved lock-on in 32 of 50 control sorties and 20 of 50 sorties with the pod.

(a) State $H_0$ and $H_1$ and carry out a two-sided two-proportion z-test at $\alpha=0.05$. [6]\newline (b) Compute a 95\% confidence interval for the reduction (unpooled standard error) and interpret it. [4]\newline (c) The same programme tested 12 different pods and this was the best result. Re-assess significance with a Bonferroni correction. [4]\newline (d) How many sorties per arm would give 80\% power to detect a reduction from 64\% to 40\%? [6]\end{prob}

\begin{soln}\textbf{(a)} $H_0: p_c=p_j$; $H_1: p_c\ne p_j$. $\hat p_c=0.64$, $\hat p_j=0.40$; pooled $\hat p=52/100=0.52$. $SE=\sqrt{0.52\times0.48\times(1/50+1/50)}=\sqrt{0.009984}=0.0999$. $z=0.24/0.0999=2.40$; two-sided $p=2[1-\Phi(2.40)]=0.016<0.05$: reject $H_0$ at the 5\% level.

\textbf{(b)} $SE_u=\sqrt{0.64(0.36)/50+0.40(0.60)/50}=\sqrt{0.009408}=0.0970$. 95\% CI $=0.24\pm1.96(0.097)=(0.050,\,0.430)$. The pod reduces lock-on probability by somewhere between 5 and 43 points: positive, but very imprecise.

\textbf{(c)} With 12 tests, the Bonferroni per-test threshold is $0.05/12=0.0042$. Since $p=0.016>0.0042$, the result is \textbf{not} significant after correction. Without correction the chance of at least one false positive among 12 null pods is $1-0.95^{12}=46\%$. The finding should be treated as a hypothesis and confirmed in a fresh, pre-registered trial.

\textbf{(d)} $n\approx\dfrac{(1.96+0.84)^2[0.64(0.36)+0.40(0.60)]}{0.24^2}=\dfrac{7.84\times0.4704}{0.0576}=64.0$, so about \textbf{64 sorties per arm} (more if a Bonferroni-corrected $\alpha$ is used).\end{soln}

\needspace{8\baselineskip}\subsection{Problem 3: Base rates and two-stage screening for insider threats \normalfont\small\color{GreyN}(20 marks)}

\begin{prob}An organisation has 50,000 accounts, of which 2 in every 10,000 are compromised on a given day. Stage 1 is a behavioural detector with sensitivity 0.95 and specificity 0.99. Accounts it flags go to an independent stage-2 detector with sensitivity 0.90 and specificity 0.98.

(a) Compute the PPV of stage 1 alone. [5]\newline (b) Compute the PPV after stage 2 and the overall sensitivity of the two-stage system. [7]\newline (c) How many accounts are flagged per day by stage 1, and by stage 2? How many true compromises are found? [5]\newline (d) Explain the operational value of the second stage in terms of analyst workload. [3]\end{prob}

\begin{soln}\textbf{(a)} $p=0.0002$. True-positive mass $0.95\times0.0002=0.00019$; false-positive mass $0.01\times0.9998=0.009998$. $\text{PPV}_1=0.00019/0.010188=\mathbf{1.9\%}$ --- 98 of every 100 alerts are false.

\textbf{(b)} The stage-2 prior is the stage-1 PPV, 0.01865. TP $=0.90\times0.01865=0.016785$; FP $=0.02\times0.98135=0.019627$. $\text{PPV}_2=0.016785/0.036412=\mathbf{46.1\%}$. Overall sensitivity $=0.95\times0.90=\mathbf{85.5\%}$ (assuming the stages' errors are independent).

\textbf{(c)} Compromised accounts per day: $50{,}000\times0.0002=10$. Stage-1 alerts: $50{,}000\times0.010188\approx509$. Stage-2 alerts: $509\times0.036412\approx18.5$. True compromises found: $10\times0.855\approx8.6$; check $18.5\times0.461\approx8.6$.

\textbf{(d)} Analysts review about 19 cases instead of about 509 per day, and roughly half of them are real rather than 2\%. This cuts alarm fatigue at the cost of missing about 1.5 compromises per day, which must be weighed against the cost of a missed insider.\end{soln}

\needspace{8\baselineskip}\subsection{Problem 4: Regression to the mean and confounding \normalfont\small\color{GreyN}(20 marks)}

\begin{prob}(a) Readiness scores have population mean 70 and test--retest reliability $\rho=0.6$. The 10\% of battalions with the lowest scores (mean 50) receive a new AI tutoring aid and score 62 on average at the next assessment. A control group of equally low-scoring battalions without the aid averages 57. Using the regression-to-the-mean model, how much improvement is expected without any intervention? Estimate the effect of the aid and comment. [10]\newline (b) Across units, simulator hours (X) correlate with exercise score (Y) at $r_{XY}=0.60$. Unit budget (Z) correlates $r_{XZ}=0.70$ and $r_{YZ}=0.75$. Compute the partial correlation $r_{XY\cdot Z}$ and interpret it. [6]\newline (c) Propose a design that would establish whether simulator hours causally improve scores. [4]\end{prob}

\begin{soln}\textbf{(a)} Expected next score with no intervention: $\mu+\rho(\bar x-\mu)=70+0.6(50-70)=58$, an improvement of 8 points from regression to the mean alone. The control group (57, +7) agrees closely with this. The treated group's +12 includes about 7--8 points of regression, so the aid's estimated effect is $62-57=\mathbf{5}$ points, not 12. Claiming 12 would have tripled the apparent benefit.

\textbf{(b)} $r_{XY\cdot Z}=\dfrac{0.60-0.70(0.75)}{\sqrt{(1-0.49)(1-0.5625)}}=\dfrac{0.075}{\sqrt{0.2231}}=\dfrac{0.075}{0.4724}=\mathbf{0.16}$. Most of the raw association reflects budget: well-funded units both use simulators more and score higher. A weak residual link remains, which may be causal or due to other confounders.

\textbf{(c)} A randomised design: within budget strata, randomly assign comparable units to additional simulator hours or to the current allocation; measure scores blind to assignment, and pre-register the outcome and analysis. Randomisation balances budget and unmeasured confounders, and the difference in mean scores then estimates the causal effect.\end{soln}

\needspace{8\baselineskip}\subsection{Problem 5: Distribution shift and evaluation design \normalfont\small\color{GreyN}(20 marks)}

\begin{prob}A soil-moisture feature was measured for 8 training sites and 8 deployment sites:\newline Training: 0.10, 0.12, 0.15, 0.18, 0.20, 0.22, 0.25, 0.30\newline Deployment: 0.20, 0.26, 0.31, 0.35, 0.38, 0.40, 0.42, 0.45

(a) Compute the two-sample Kolmogorov--Smirnov statistic $D=\max_x|F_{\text{train}}(x)-F_{\text{deploy}}(x)|$. [6]\newline (b) Using the 5\% critical value $1.36\sqrt{(n+m)/(nm)}$, decide whether the distributions differ. [3]\newline (c) A deep model scores 97\% on training-domain test data. Design a pre-deployment evaluation protocol of at least five steps that would catch an Afghan-to-Iraq style failure, and state what result would block deployment. [11]\end{prob}

\begin{soln}\textbf{(a)} Evaluate both empirical CDFs at each data point. At $x=0.25$: $F_t=7/8=0.875$, $F_d=1/8=0.125$, difference 0.75. At $x=0.30$: $F_t=1.000$, $F_d=2/8=0.25$, difference 0.75. All other points give smaller differences (e.g. $x=0.22$: 0.625). $\mathbf{D=0.75}$.

\textbf{(b)} Critical value $=1.36\sqrt{16/64}=0.68$. Since $D=0.75>0.68$, reject equality at the 5\% level: the deployment sites are systematically wetter, so the training data do not represent the deployment domain.

\textbf{(c)} A model answer covers: (1) collect a labelled sample from the deployment domain (at least 100--200 cases, all expected terrains, seasons and threat variants) before fielding; (2) run KS or similar shift tests on every key feature, with a pre-set threshold; (3) report per-class recall on the deployment sample, especially the threat class, with confidence intervals, not only overall accuracy; (4) compare against a simple baseline and the training-domain score, and treat a drop of more than a set margin (e.g. 5 points) as a failure; (5) stress-test with adversarial and degraded inputs (sensor noise, camouflage); (6) pilot in advisory mode with human confirmation and log every disagreement; (7) re-train or apply domain adaptation if a test fails. \textbf{Blocking rule:} deployment-domain threat recall below the operational requirement, or a significant shift on safety-critical features without a successful domain-adaptation re-test, halts fielding.\end{soln}

\end{document}